\documentclass[11pt]{article}
\usepackage{finprobs}

\title{Theoretical Improvements to Unified Rough/Multifractal\\ Volatility Models: A Synthesis}
\author{finprobs Research Project}
\date{July 2026}

\begin{document}
\maketitle

\begin{abstract}
The log S-fBM model of \citet{wu2022rough} unifies rough volatility (fractional-Brownian-driven log-volatility, Hurst $H>0$) and the Multifractal Random Walk / log-normal cascade ($H\to0$) as limits of a single Gaussian process. Unification alone is not a theoretical advance: it does not by itself explain why $H$ takes the values it does, whether the model's multifractality is a genuine temporal phenomenon, or whether it improves anything practitioners calibrate against. This paper synthesizes a four-part investigation (detailed in companion papers in this series) into exactly these questions, conducted by an orchestrated multi-agent research process with adversarial citation auditing and direct empirical/computational verification. We correct a persistent mischaracterization of the model's own mathematical structure, report a decisive asymmetry in how well-supported its two dimensions are, and state two small but genuine new theoretical results not found anywhere in the literature searched.
\end{abstract}

\section{Background and motivation}

Rough volatility \citep{gatheral2018volatility,eleuch2018microstructural} and multifractal cascade models \citep{bacry2003multifractal} have historically been treated as competing descriptions of the same stylized facts: volatility clustering, long memory, fat tails, and apparent multiscaling. \citet{wu2022rough} showed these are not competitors but two limits of a single one-parameter family -- a mathematically clean result. The question this series addresses is what, beyond that unification, is actually gained: does it explain the empirically observed values of $H$, is the multifractality it produces genuinely temporal, can its parameters be derived rather than fit, and does it solve any practical calibration problem that the separate models could not.

We organize the investigation into four directions, each the subject of a companion paper in this series:
\begin{enumerate}
\item[\textbf{P1}] Distributional versus temporal multifractality.
\item[\textbf{P2}] Scale-dependence of $H$ as a renormalization-group-like flow across aggregation.
\item[\textbf{P3}] Microstructural derivation of the roughness and intermittency parameters.
\item[\textbf{P4}] Empirical/pricing payoff for joint SPX/VIX calibration.
\end{enumerate}

\section{Methodology}

The investigation was conducted by approximately 26 independent large-language-model research and audit agents across five rounds: (i) eight diverse-exploration agents, two independent angles per direction, deliberately kept unaware of each other's approach to preserve genuine independence; (ii) manual synthesis into an approach registry grouped by underlying mechanism, used to redirect effort away from premature convergence; (iii) ten deepening and dedicated adversarial-audit agents, checking citation accuracy, conflation of distributional and temporal multifractality, circularity in derivations, and cross-direction consistency; (iv) four cross-pollination agents connecting compatible threads across directions; (v) eight follow-up agents on problems the first four rounds newly opened. Every citation flagged as surprising or load-bearing was independently re-verified against its primary source (arXiv abstracts and full text fetched directly, not trusted from any single agent's summary); two accuracy errors -- an author misattribution and a mischaracterized maturity range in one cited result -- were caught this way and corrected before publication. A subset of testable claims was additionally verified by direct computation: closed-form numerical checks, Monte Carlo simulation against known ground truth, and analysis of a freshly collected fifteen-year real-market dataset (S\&P 500 index and 25 large-cap single names) that no cited paper uses.

\section{A necessary correction: what the model actually is}

The informal description of log S-fBM as a ``Box--Cox/Tsallis-style deformation'' -- used in early framings of this investigation and common in secondary discussion of the model -- is not accurate. Reading \citet{wu2022rough} directly gives log-volatility $\omega_{H,T}(t)$ as a \emph{stationary Gaussian process}, with covariance
\begin{equation}
C_\omega(\tau) = \frac{\nu^2}{2}\left[T^{2H}-\tau^{2H}\right], \qquad |\tau|<T, \qquad \nu^2 = \frac{\lambda^2}{H(1-2H)}, \label{eq:cov}
\end{equation}
and volatility measure $M_{H,T}(dt)=e^{\omega_{H,T}(t)}dt$. There is no separate nonlinear marginal deformation function; $H$, $\lambda^2$, and the decorrelation scale $T$ are three structurally independent free parameters (mildly anti-correlated only empirically, not by theoretical necessity). As $T\to\infty$, the process converges to standard fBM (\citealp{wu2022rough}, Prop.~1); as $H\to0$ with $\lambda^2$ fixed, \eqref{eq:cov} converges to the exactly logarithmic covariance $\lambda^2\log(T/\tau)$ (\citealp{wu2022rough}, Prop.~2) -- the genuine mechanism generating multifractality via Gaussian multiplicative chaos \citep{kahane1985chaos}, not a curvature parameter on a marginal distribution. The source paper is purely a probability/statistics paper (GMM estimation applied to realized volatility across 24 indices and 296 stocks); it contains no options-pricing content, and its own authors flag the aggregation question addressed in P2 as explicit future work, confirming it as a genuine open problem rather than one this investigation invented.

\section{Synthesis of findings}

\subsection{The central asymmetry}

Combining all four companion papers, one conclusion is decisive and was independently corroborated from three unrelated angles:

\begin{finding}[Central asymmetry]
The roughness ($H$) dimension of log S-fBM is moderately supported by both theory and data: a rigorous microstructural derivation exists at the pure-fBM endpoint (P3, \citealp{eleuch2018microstructural}), and a rigorous, now simulation-confirmed aggregation mechanism explains why index volatility is smoother than any single name, conditional on an empirical input (P2, \citealp{zarhali2025nested}). The multifractal/intermittency ($\lambda^2$) dimension is the weakest-supported part of the entire framework, undermined simultaneously and independently: it is not established whether the model's own multiscaling is genuinely temporal rather than a fat-tailed-marginal artifact (P1); no microstructural mechanism, Hawkes-based or otherwise, has been found to produce it, and two specific candidates were proven not to work (P3); and a first-principles derivation shows $\lambda^2$ is calibration-inert for the one practical payoff most likely to justify it (P4).
\end{finding}

Two independent derivation attempts -- P2's attempt to derive why the common factor is smoother than idiosyncratic volatility, and P3's attempt to derive a multiplicative Hawkes-hierarchy mechanism for the cascade -- hit structurally the same wall: both had to posit exogenous structure to obtain the desired result rather than deriving it from first principles. This recurrence across independent efforts suggests a genuinely hard, not merely under-explored, class of problem.

\subsection{Novel results}

Two small but genuine theoretical results emerged that are not present anywhere in the literature searched, both now verified by direct computation rather than resting on derivation alone.

\begin{theorem}[P4]
At fixed $H$, the intermittency parameter $\lambda^2$ cancels out exactly of the ratio between short-maturity and VIX-horizon log-volatility covariance under \eqref{eq:cov}; the correlation-timescale parameter $T$ governs that ratio instead.
\end{theorem}

\begin{theorem}[P3]
Linear superposition of independent nearly-unstable Hawkes-process scaling limits remains jointly Gaussian under time-invariant mixing weights, and therefore cannot produce genuine multifractal curvature in the structure function of $\log\sigma_t$ -- narrowing the open problem of a microstructural cascade derivation to one precisely specified missing double-scaling limit theorem.
\end{theorem}

\subsection{Cross-direction consistency}

An adversarial audit run across all four directions together identified a specific, reportable tension: P2's use of single-stock $H\approx0.01$ as evidence of genuine ``super-roughness'' sits inside exactly the regime P1 found to be dominated by distributional (not temporal) confound in the closely related rough Bergomi model. If this transfers to log S-fBM, P2's empirical anchor may be partly an estimation artifact rather than pure dynamics -- a dependency P2 does not resolve and flags explicitly. A second tension, discovered by direct cross-paper derivation, is documented in P4: the same aggregation mechanism that explains P2's central claim likely erodes, rather than reinforces, the $T$-lever proposed in P4 as the practical payoff of unification. Finally, the \citet{gatheral2018volatility} $H\approx0.1$ estimate for the S\&P~500 is used as an empirical anchor in both P2 and P4; it should be read as one data point in two contexts, not two independent corroborations.

\section{Findings summary}

\begin{itemize}
\item \labelestablished. The model's actual mathematical structure (Section~3), correcting prior informal descriptions.
\item \labelestablished. The central asymmetry between the $H$ and $\lambda^2$ dimensions (Finding~1), corroborated across three independent lines of investigation.
\item \labelestablished. Both novel theoretical results (Theorems~1 and~2), each verified by direct computation.
\item \labelplausible. The specific cross-direction tensions identified in Section~4.3, each with a stated, checkable follow-up.
\end{itemize}

\section{Future directions}

Ranked by how directly they would move the open questions forward, given what is now known across all four companion papers:
\begin{enumerate}
\item Run the boundary-corrected genuineness diagnostic proposed in P1 properly (full Self--Liang likelihood-ratio test, high-frequency realized variance, larger sample) rather than the simplified version validated here.
\item Attempt the joint double-scaling Hawkes limit theorem specified in P3 -- the missing object for a positive microstructural derivation of the cascade parameters is now precisely stated.
\item Build a minimal log S-fBM option-pricing engine and test the $T$-lever hypothesis from P4 directly against public SPX/VIX data, with an explicit parameter-matched overfitting control.
\item Test the branching-ratio prediction from P2's speculative economic mechanism, and the $T_{\text{index}}$ prediction from P4's cross-direction tension, both of which are new and currently untested.
\item Reconcile P1's model-level finding (classical MRW is temporal-dominant) with the market-level finding that observed multifractality is mainly distributional -- either by testing whether log S-fBM with realistic parameters reproduces the latter, or by identifying the methodological source of the discrepancy.
\end{enumerate}

\section*{Acknowledgments}
This paper and its companions were produced by an orchestrated multi-agent AI research investigation. All citations were independently re-verified against primary sources; all numerical claims were computed, simulated, or estimated directly and confirmed before inclusion, with limitations reported honestly rather than omitted, per the investigation's own stated reporting standard.

\bibliographystyle{plainnat}
\begin{thebibliography}{99}

\bibitem[Bacry and Muzy(2003)]{bacry2003multifractal}
E.~Bacry and J.-F.~Muzy.
\newblock Log-infinitely divisible multifractal processes.
\newblock \emph{Communications in Mathematical Physics}, 236(3):449--475, 2003.

\bibitem[El~Euch et al.(2018)]{eleuch2018microstructural}
O.~El~Euch, M.~Fukasawa, and M.~Rosenbaum.
\newblock The microstructural foundations of leverage effect and rough volatility.
\newblock \emph{Finance and Stochastics}, 22(2):241--280, 2018.

\bibitem[Gatheral et al.(2018)]{gatheral2018volatility}
J.~Gatheral, T.~Jaisson, and M.~Rosenbaum.
\newblock Volatility is rough.
\newblock \emph{Quantitative Finance}, 18(6):933--949, 2018.

\bibitem[Kahane(1985)]{kahane1985chaos}
J.-P.~Kahane.
\newblock Sur le chaos multiplicatif.
\newblock \emph{Annales de la Facult\'e des Sciences de Toulouse}, 9(2):105--150, 1985.

\bibitem[Wu et al.(2022)]{wu2022rough}
Q.~Wu, J.-F.~Muzy, and E.~Bacry.
\newblock From rough to multifractal volatility: the log S-fBM model.
\newblock \emph{Physica A}, 604:127919, 2022. arXiv:2201.09516.

\bibitem[Zarhali et al.(2025)]{zarhali2025nested}
O.~Zarhali, C.~Aubrun, E.~Bacry, J.-P.~Bouchaud, and J.-F.~Muzy.
\newblock A nested factor model for equity markets: reconciling multifractal stock returns and rough index volatilities.
\newblock \emph{arXiv:2505.02678}, 2025.

\end{thebibliography}

\end{document}
